% ============================================================================
% CH — GROUNDING THE ITEM SIDE: fI-ProtoMF              [slug: fi-protomf]
% STATUS: theory/design sections material-rich (dc01 hardened through SC.5);
% results sections BLOCKED on S0.7 fleet + SC.6+ runs.
% Lineage step 1-4 of the committed outline (2026-07-11_2351). Presentation
% order per user: present fully -> why -> demonstrate -> route-1 lens.
% DELTA PRINCIPLE (Matteo, 2026-07-14): this is the FULL-BUILD chapter — the
% one place where what/why/how is developed exactly and precisely. Every
% later lineage chapter (fU, merge) teaches only its DELTA against this one
% and shrinks accordingly (page goals descend 10-12 -> 8-10 -> 6-8).
% HEADERS REWORKED 2026-07-14 (granularity: headers follow mass): 8 -> 6
% sections; four-component inventory folded into The Model as a block;
% route-1 lens folded into Explanations as the closing bridge block.
% Material map: ../outline.md#fi-protomf
% ============================================================================
\cleardoublepage
\chapter{Grounding the Item Side: fI-ProtoMF (10--12)}   % page goal — scaffolding, DELETE before submission
\label{ch:fi}

\section{Why the Item Side First: Taste Is Revealed, Not Stated}
\label{sec:fi-why}
\vault{2026-07-11_2350_interactions-explained-by-item-features-taste-revealed}
% MOVED to open the chapter (Matteo, 2026-07-14): motivate before building.
% The core argument: interactions are explained by item features; user taste
% must be discovered from behavior; demographics are a too-loose proxy.
% LOAD-BEARING BEYOND THIS CHAPTER: the fU chapter's opening calls back to
% this section to justify itself (users composed from revealed behavior —
% histories — not stated attributes; also why afU enters last and weakest in
% the merge chapter). Write it to carry that weight.

\section{The Model: Feature-Composed Item Factors}
\label{sec:fi-model}
\label{sec:fi-inventory}   % old inventory section parked here (block 3)
\vault{2026-06-16_1715_feature-injection-design-fork, 2026-07-11_1526_facet-a-claim-boundary-coldstart-2x2}
% Content blocks:
%  1. Construction: q_i = sum of feature-word embeddings e_f (+ ID row e_ID)
%     on the I-ProtoMF host; keystone reduction f*(F=0,+ID) == host
%     (bit-identity); ablation arms ids / noid / f0.
%  2. Terminology home (bullet below).
%  3. The four-component inventory — what is grounded, what is not:
%     t* grounded, u-hat derived-legible, t-hat attributable-but-ungrounded-
%     units, u* ungrounded — immediate transparency about the half-opaque
%     part; score-share disclosure (Block-I/Block-U reporting rule).
\begin{itemize}
    \item Terminology home: attribute = human-readable catalogue property; feature = its encoded, model-facing row. \vault{2026-07-13_2113_why-f-not-attribute-in-model-names-mechanism-vs-contribution}
    \item ONE sentence (DIRECTIVE) on why the model names keep the \emph{f}: choosing attribute rows is what buys interpretability. \vault{2026-07-13_2113_why-f-not-attribute-in-model-names-mechanism-vs-contribution}
\end{itemize}

\section{Prototypes as Re-Coordinatization}
\label{sec:fi-recoordinatization}
\vault{2026-07-11_2214_thesis-recoordinatization-narrative-and-3d-figure, 2026-07-11_2201_latent-directions-meaning-prototype-recoordinatization}
% "Axes meaningless, directions meaningful"; behavior fills the space ->
% prototypes give addresses -> features give names. Committed d=3 emergence
% figure + caption argument (why not post-hoc naming of latent directions).

\section{Reading the Model: Attribution and Its Limits}
\label{sec:fi-readouts}
\vault{2026-07-11_2132_dc01-id-residual-three-roles, 2026-07-11_2310_dc01-share-identifiability-problem-and-thesis-framing}
% Per-feature shares c_{f,k}; correlated features -> share identifiability
% (grouped shares, determinacy annotations — full treatment in the
% hidden-effects chapter); the ID residual's three roles (honesty meter /
% popularity reading / mis-tag probe); hard boundary: ID rows never enter
% prototype profiles.

\section{Results}
\label{sec:fi-results}
% [PENDING S0.7 + SC.6+] accuracy vs fleet; ids/noid/f0 ablations; cold-item
% performance; feature-explained fraction. Scrutiny numbers are working
% evidence only — thesis rows come from dedicated runs.

\section{Explanations in Practice}
\label{sec:fi-explanations}
\label{sec:fi-route1}   % old route-1 lens section parked here (closing block)
\vault{2026-06-08_1149_prototype-naming-and-similarity-routes}
% DELTA ONLY (2026-07-14): the explanation system + host baseline are
% demonstrated in the eval-framework chapter — here show only the EXTENDED
% capability (parameter-derived item-prototype names), no host re-demo.
% Content blocks:
%  1. The explanation demonstrated as it stands, opaque user component shown
%     honestly; fashion showcase (hm_3_month qualitative).
%  2. CLOSING BRIDGE — a lens ahead: user prototypes through the tie
%     (route-1: projecting user prototypes into named item-prototype space;
%     enabled-not-enforced; extrapolation caveat — projection trained on
%     embeddings, not prototypes). Motivates the user side -> next chapter.
